% DELTA23: quadratic orbital identity; root independently certified all13 paths.
\subsection{The harmonic orbital identity and its shared periodic braid}

\paragraph{D23.C18 (the quadratic orbital cofaces and residual).}
Retain the literal four-strand marking
\[
 a=\sigma_1,\quad b_0=\sigma_2,\quad c=\sigma_3,\qquad
 x=a^2,\ y=b_0^2,\ K=c^2,\ Z=(ab_0)^3,\ V=(b_0c)^3,
\]
and put $s=ab_0a$, $r=ab_0$, $z_4=(ab_0c)^4$.
There is a quadratic orbital map
\[
 \mathcal O^H_n:W_{n,2}\longrightarrow M_{0,n+2},\qquad
 (z_1,\ldots,z_n)\longmapsto(-z_1,0,z_1,z_2,\ldots,z_n).
\tag{D23.C18.1}
\]
Use the positive source chamber, send the last target point to
infinity by $p\mapsto1/(p-z_3)$, and straighten the remaining
strands in reverse label order. Unlike D23.C16, there is no
final common conjugation. All initial target points are real,
and their transformed order is uniformly $4<3<2<1$ throughout
the positive chamber. Thus all face connecting paths have the
same marking without additional conjugators.

For the cyclic source basis
\[
 U=x_{01}^2,\qquad Y_0=x_{12},\qquad
 Y_1=x_{01}x_{12}x_{01}^{-1},
\]
the three middle faces have the following literal $B_4$ images:
\[
\begin{array}{c|ccc}
 &U&Y_0&Y_1\\\hline
 \widetilde\partial_1&z_4&Zz_4^{-1}&1\\
 \widetilde\partial_2&z_4&xK^{-1}z_4^{-1}&Vz_4^{-1}\\
 \widetilde\partial_3&Z&xZ^{-1}&yZ^{-1}.
\end{array}
\tag{D23.C18.2}
\]
The exceptional zeroth face sends its two generators to
$K,Zz_4^{-1}$. The fourth face sends them to $s,K$.
Here the partial half rotation of $z_1$ exchanges the two
outer points of the cluster $(-z_1,0,z_1)$ and gives $s$;
it satisfies $s^2=Z$.

Let $h=(1,f)$ be an actual multiplier-one full-GT element and put
\[
 p=\Psi_2(h),\qquad b=\Psi_3(h)/2,\qquad
 J_f=f(s,x)\in F,\qquad B^H_f=x^{-2p}J_f\in F.
\tag{D23.C18.3}
\]
Then $(J_f)^{\rm ab}=(p,-p)$ and
$(B^H_f)^{\rm ab}=(-p,-p)$. The harmonic orbital identity is
\[
 \boxed{\quad
 f(K,Z)B^H_f(x,V)K^{-p}f(s,K)
                  =B^H_f(x,y)Z^p z_4^{-p}.
 \quad}
\tag{D23.C18.4}
\]
It is a literal identity in $\widehat B_4$ and therefore also
holds after imposing $z_4=1$. No assertion that $B^H_f$ is a
frame of an actual gentle descent is needed or made.

Let $\beta_h$ be the original pure harmonic frame and define
\[
 D^H_h=\beta_h(x,y)^{-1}B^H_f(x,y),\qquad g_H=f(s,K).
\]
In the sphere quotient one has
\[
 \boxed{\quad
 (sD^H_hg_H^{-1}c g_HD^{H,-1}_h)^3=1.
 \quad}
\tag{D23.C18.5}
\]
For the original signed modular quotient $q_3$, its residual is
\[
 q_3(D^H_h)\in\langle S\rangle\simeq C_4,
 \qquad
 \epsilon_2(q_3(D^H_h))=a_2(h)+b=d_2(h).
\tag{D23.C18.6}
\]
Thus C17 and C18 give correlated five-point periodic
factorizations with residual parities $e_3$ and $d_2$.
They do not by themselves prove equality of these parities.

\begin{proof}
Use the same coface paths as in D23.C16, now with degree two.
The three lifted cofaces are
$(az_1,z_1,z_2)$, $(z_1,cz_1,z_2)$ and
$(z_1,z_2,cz_2)$, with real splitting constants $a$ small and
$c>1$ close to one. In this sentence the two constants are not
braid generators. The positive-root lollipop has the same
boundary twist description as in C16. The only other root
is $-z_1$; the fixed point $0$ remains marked. Consequently
the first two rows in (D23.C18.2) are the stated boundary
twists, and the third follows from $sxs^{-1}=y$ and $s^2=Z$.
The zeroth face follows the positive collision arcs, and the
fourth is the half rotation followed by the same collision.

There is an independent exact geometric certificate of all
thirteen generator paths. The program
\path{code/certify_harmonic_orbital_paths.py} uses the cubic
certificate's rational Sturm engine, with six edges of the
regular twelve-gon for a half rotation and all twelve for $U$.
It uses the target coordinates $(-z_1,0,z_1,z_2,z_3)$ and the
projection factor $1+i/10$. The positive collision diamonds
and positive-chamber connecting paths have the same elementary
separation properties as in C16. Exact rational isolation of
every crossing polynomial, its derivative sign, and its
imaginary-coordinate sign proves all thirteen literal crossing
words, with no sampled inference. The full certificate is
\path{records/harmonic_orbital_paths_exact.json}.
The separate faithful Artin replay
\path{agents/characters/code/check_harmonic_orbit_coface_artin.py}
checks every simplification in (D23.C18.2), with its exact
central exponent; it also checks all braid identities used below.

In the index-two Schreier abelianization the coefficient of $U$
is zero and those of $Y_0,Y_1$ are $(p,-p)$.
The filling $U\mapsto1,Y_0\mapsto x,Y_1\mapsto y$ sends $f$
to $J_f=f(s,x)$, since $s^2=1$ and $sxs=y$ in the central
three-strand quotient. This gives the asserted abelianization.
The equality lifts to the marked free pure subgroup of $B_3$:
both sides have total braid exponent zero, so their possible
central discrepancy is zero. Equivalently the $U$ coefficient
is zero and $s^2=Z$ is central in the literal lift.

Rows one, two and three therefore give, respectively,
\[
 Z^p z_4^{-p},\qquad J_f(x,V)K^{-p},\qquad J_f(x,y).
\]
The two central exponents in row two cancel because
$p+(-p)=0$, as do the $Z$ exponents in row three.
The commutator-word evaluation in the zeroth face is
$f(K,Zz_4^{-1})=f(K,Z)$.
Apply Enriquez's lifted pentagon (51) in the form
$\partial_0(f)\widetilde\partial_2(f)\partial_4(f)
=\widetilde\partial_3(f)\widetilde\partial_1(f)$.
Then substitute $J_f=x^{2p}B^H_f$ and cancel the common
factor $x^{2p}$, which commutes with $K,Z$ and therefore
with $f(K,Z)$. This proves (D23.C18.4).
Its two sides have the same literal braid exponent $-10p$.

Put $R=f(K,Z)$ and $M_H=B^H_f(x,V)K^{-p}$.
Both $x$ and $V$ commute with $c$, so $M_H$ commutes with $c$.
Equation (D23.C18.4) gives, in the sphere quotient,
\[
       (B^H_f)^{-1}R=Z^p g_H^{-1}M_H^{-1}.
\]
The standard full-GT braid action has
$\Phi(s)=\beta_hs\beta_h^{-1}$ and
$\Phi(c)=RcR^{-1}$. The first equality lifts literally
from the three-strand central quotient by braid exponents;
the second is the standard Drinfeld braid formula, using
the two-cycle to reverse $f(Z,K)$.
Since $(sc)^3=z_4$, and $Z$ centralizes $s,D^H_h$, it follows
that
\[
 \Phi(sc)=\beta_h Z^p
  (sD^H_hg_H^{-1}c g_HD^{H,-1}_h)Z^{-p}\beta_h^{-1}.
\]
Cubing proves (D23.C18.5).

It remains to identify the finite residual. Write
$q_3(x)=U^2$ and $q_3(s)=S$, allowing the central lift sign
of $s$, which does not change a commutator-word evaluation.
In the projective affine realization over $\mathbb F_4$ used
in A.87 and A.106,
\[
 x=(0,\omega),\qquad s=(c_0,1),\qquad
 c_0=(1-\omega)^{-1}.
\]
A.90 and the two-cycle give
$J_f=f(s,x)=(c_0b,1)$ in this affine quotient (the inverse
has the same translation in characteristic two).
The inverse image of the subgroup generated by this projective
$s$ in $\operatorname{SL}_2(\mathbb F_3)$ is
$\langle S\rangle\simeq C_4$. Hence
\[
 W_H:=q_3(B^H_f)^{-1}U^{2p}
     =q_3(J_f)^{-1}\in\langle S\rangle,
       \qquad \epsilon_2(W_H)=b.
\]
The equality uses $U^3=1$:
$q_3(B^H_f)^{-1}U^{2p}=q_3(J_f)^{-1}U^{6p}$.
For the original frame
$W_h=q_3(\beta_h)^{-1}U^{2\rho(h)}\in\langle S\rangle$.
A.102 supplies the full-GT congruence $\rho(h)=p\pmod3$.
Consequently
\[
 q_3(D^H_h)
   =W_hU^{-2\rho(h)-4p}q_3(J_f)
   =W_h W_H^{-1}.
\]
Taking the quotient $C_4\to C_2$ proves (D23.C18.6).
This finite computation does not presume an automorphism of
$\widehat B$ extending the filled index-three subgroup action.
\end{proof}

\paragraph{D23.C19 (the shared periodic conjugacy).}
Let $P=(ab_0c)^2$. There is the literal braid identity
\[
                P(sc)P^{-1}=rK.
\tag{D23.C19.1}
\]
Write
\[
 H_h=sD^H_hg_H^{-1}c g_HD^{H,-1}_h,\qquad
 E_h=rD_hg^{-1}KgD_h^{-1},
\]
with the cubic notation $D_h,g,A_h,b$ of C17.
Then both residual periodic elements satisfy the specific
common conjugacy
\[
 E_h=L_hH_hL_h^{-1},\qquad
 L_h=Z^{-b}A_h^{-1}\Phi(P)\beta_hZ^p.
\tag{D23.C19.2}
\]
Indeed C17 and C18 express $\Phi(rK)$ and $\Phi(sc)$ as
the stated conjugates of $E_h,H_h$; applying $\Phi$ to
(D23.C19.1) gives (D23.C19.2). The braid identity follows
also from conjugation by $ab_0c$ shifting $a$ to $b_0$ and
$b_0$ to $c$: $sc=(ab_0c)a$ and $rK=(ab_0c)c$.
Thus the two residual torsion conditions are coupled by the
same original action. Reading their two frame parities from
this common conjugacy, or from an equivalent finite quotient,
remains the unresolved comparison; arbitrary unrelated
conjugators would discard precisely this extra data.

\paragraph{D23.C20 (why the affine shared conjugacy alone is insufficient).}
The shared conjugacy C19 imposes no parity constraint after applying
$\eta_3$ and retaining only the supplied full-GT covariance
$\eta_3\Phi=\operatorname{Ad}(0,U^{2\rho})\eta_3$.
This is stronger than the separate observation that both cube
tests admit their two possible defect signs.

Indeed suppress $\eta_3$ from the notation and put
\[
 V_0=(0,U^{2\rho}),\quad
 T=A_h^{-1}V_0,\quad W=\beta_h^{-1}V_0,\quad z=Z,
 \quad d_E=rK,\ d_H=sc,\ p_0=P.
\]
Here $T,W,z$ are linear affine matrices, and $z=-I$ commutes
with both $T,W$. C17--C19 and the covariance give exactly
\[
 E_h=\operatorname{Ad}(z^{-b}T)(d_E),\qquad
 H_h=\operatorname{Ad}(z^{-p}W)(d_H),\qquad
 L_h=(z^{-b}T)p_0(z^{-p}W)^{-1}.
\tag{D23.C20.1}
\]
For example $A_h^{-1}V_0=T$ converts the conjugator of
$\Phi(rK)$ directly into the first expression; the harmonic
expression is identical with $\beta_h,W,p$ in place of
$A_h,T,b$. These formulas use no equality between the two
frame parities. Since $p_0d_Hp_0^{-1}=d_E$, substitution proves
$E_h=L_hH_hL_h^{-1}$ identically for every
$T\in C_6$, $W\in C_4$ and every $b,p$.
Thus the affine translation class in this particular conjugacy
cannot recover the missing parity by itself. Additional
same-word constraints on the full orbital factors have not
been discarded by this statement and remain a separate test.
